\begin{table}[htbp]
\centering
\small
\caption{Bairros com maior incidência média anual de dengue, 2013--2025.}
\label{tab:incidencia}
\begin{adjustbox}{max width=\linewidth}
\begin{tabular}{rlrrr}
\toprule
& Bairro & População (2022) & Casos & Incidência média anual/100 mil \\
\midrule
1 & Campina do Barreto & 8.452 & 1.091 & 992,9 \\
2 & Estância & 7.496 & 643 & 659,9 \\
3 & Ibura & 46.679 & 3.388 & 558,3 \\
4 & Cohab & 61.408 & 4.345 & 544,3 \\
5 & Bomba do Hemetério & 8.415 & 583 & 532,9 \\
6 & Coelhos & 5.731 & 396 & 531,5 \\
7 & Morro da Conceição & 8.209 & 566 & 530,4 \\
8 & Afogados & 31.901 & 2.107 & 508,1 \\
9 & Areias & 26.537 & 1.746 & 506,1 \\
10 & Beberibe & 7.047 & 448 & 489,0 \\
11 & Arruda & 12.988 & 822 & 486,8 \\
12 & Cajueiro & 5.106 & 318 & 479,1 \\
13 & Água Fria & 38.086 & 2.208 & 446,0 \\
14 & Jordão & 19.808 & 1.126 & 437,3 \\
15 & Joana Bezerra & 12.495 & 678 & 417,4 \\
\bottomrule
\end{tabular}
\end{adjustbox}
\par\smallskip\begin{minipage}{0.95\linewidth}\footnotesize Casos confirmados com bairro identificado (inclui casos localizados pelo centroide). Apenas bairros com pelo menos 5.000 habitantes. População estimada por ponderação areal dos setores censitários.\end{minipage}
\end{table}
